% Part: second-order-logic
% Chapter: metatheory
% Section: introduction

\documentclass[../../../include/open-logic-section]{subfiles}

\begin{document}

\olfileid{sol}{met}{int}

\olsection{Pambuka}

Logika orde kapisan nduweni sawetara sipat kang becik. Kita
ngerti yen logika iki ora bisa diputusake, nanging paling ora
bisa diaksiomatisasi. Tegese, ana sistem bukti kanggo logika
orde kapisan kang sahih lan jangkep, yaiku ngasilake relasi
!!a{derivability}~$\Proves$ kanthi sipat: kanggo saben
himpunan !!{sentence}s~$\Gamma$ lan !!{sentence}~$!A$,
$\Gamma \Entails !A$ yen lan mung yen
$\Gamma \Proves !A$. Iki mligine ateges validitas logika
orde kapisan iku !!{computably
  enumerable}. Ana fungsi komputabel~$f\colon \Nat \to
\Sent[L]$ kang nilai-nilai~$f$ persis kabeh !!{sentence}s valid
saka~$\Lang{L}$. Iki amarga !!{derivation}s bisa
dienumerasi, lan kang !!{derive} siji~!!{sentence}
sabanjure dipetakake marang !!{sentence} kasebut.
Logika orde kapindho nduweni daya ungkap luwih gedhe tinimbang
logika orde kapisan, mula umume luwih ruwet kanggo nemtokake
validitase. Sejatine, kita bakal nuduhake yen logika orde
kapindho ora mung ora bisa diputusake, nanging validitase
malah ora !!{computably
  enumerable}. Iki ateges ora ana sistem bukti kang sahih
lan jangkep kanggo logika orde kapindho (sanadyan sistem
bukti kang sahih nanging ora jangkep ana, lan pancen dadi
objek panliten kang wigati).

Logika orde kapisan uga nduweni rong sipat liyane: logika iki
kompak (yen saben subhimpunan winates saka himpunan~$\Gamma$
kang ngemot !!{sentence}s bisa dicukupi, mula $\Gamma$
dhewe bisa dicukupi), lan Teorema L\"owenheim--Skolem
uga bener kanggo logika iki (yen $\Gamma$ nduweni model tanpa wates,
mula uga nduweni model !!a{denumerable}). Loro asil iku
ora bener kanggo logika orde kapindho. Sebabe, logika
orde kapindho bisa nyatakake fakta ngenani ukuran
!!{domain}s kang ora bisa dinyatakake dening logika orde
kapisan.

\end{document}
